% SPDX-License-Identifier: CC-BY-4.0
% =====================================================================
% Section 8 of D. Joyner and Claude, "A hitch-hiker's guide to the
% metaplectic trace formula for SL(2): basic theory" (November 2026).
% New in the style-guide pass (style guide, Section 2.7). It restates
% results of Sections 2-7 with the status given there and lists only open
% questions that Sections 1-7 already name. It is new text and is listed
% for refereeing in ../revision_log.md (Stage 8).
% =====================================================================

\section{Summary of results and open questions}\label{sec:summary}

This section collects the main results of the paper, with the status
that the text gives them, and the questions that the paper leaves open.
The status of each statement in detail is in the subsections ``Status of
the results'' at the end of Sections~\ref{sec:background}--\ref{sec:stable}.

\subsection*{Results}

\begin{enumerate}[label=\arabic*.]
\item \emph{The cover} (Section~\ref{sec:background}). The double cover
  of $G(\A)$ is built from Kubota's cocycle, with its topology fixed
  intrinsically, and it splits over $G(\Q)$ (Lemmas~\ref{lem:kubota}
  and~\ref{lem:global}). The cocycle and the existence of the topological
  cover are taken from~\cite{Ku67,Mo68}. The rest is proved in the text.
\item \emph{Rank $0$} (Section~\ref{sec:rank0}). For the cover of a torus
  the trace formula is Poisson summation, twisted by a genuine
  automorphic character (Theorem~\ref{thm:tf-rank0} and
  Lemma~\ref{lem:transfer}). For the split torus Weil's character
  $\chi_\psi$ gives a canonical choice (Lemma~\ref{lem:weil}, whose
  parts~(a) and~(b) are Weil's).
\item \emph{The geometric side} (Section~\ref{sec:geometric}). The coarse
  geometric expansion is proved directly from the height inequality, with
  truncation parameter $T\ge0$ (Lemma~\ref{lem:height-ineq} and
  Theorem~\ref{thm:coarse}). Each class contributes a polynomial in $T$ of
  degree at most one, written down in
  Propositions~\ref{prop:hyperbolic}, \ref{prop:elliptic}
  and~\ref{prop:central}. The fine expansion, Theorem~\ref{thm:fine}, is
  quoted from~\cite{Ar86,Li14a}. Its part~(a) is derived again in
  Proposition~\ref{prop:unipotent} when the ring of $S_k$-integers is
  principal.
\item \emph{The spectral side} (Section~\ref{sec:spectral}). The
  characters fixed by the Weyl group are the Weil characters
  $\chi_{\psi_b}$ (Lemma~\ref{lem:w-action}). The normalizing factor of the
  intertwining operator is $L^S(2\lambda,\chi^2)/L^S(2\lambda+1,\chi^2)$
  (Lemma~\ref{lem:intertwining}). The residual spectrum is the sum of the
  even Weil representations $r^+_{\psi_{-b}}$ (Theorem~\ref{thm:residual}).
  Parts~(a) and~(b) of that theorem, and the matching of parameters
  in~(c), are proved here. The rest of~(c) is taken
  from~\cite{Ge76,GPS80,Wa80}. The trace formula is
  Theorem~\ref{thm:tf}. The spectral decomposition, the Eisenstein series
  and the polynomial behavior of the spectral terms are taken
  from~\cite{La76,MW95,Ar80,Li13}.
\item \emph{The invariant form} (Section~\ref{sec:invariant}). The
  invariant distributions are built by Fourier inversion on the abelian
  cover of $A$ (Theorem~\ref{thm:IA}). The invariant trace formula is
  Theorem~\ref{thm:itf}, with Corollary~\ref{cor:itf-final}, and
  $M(w,0)=-1$ for every character fixed by the Weyl group
  (Proposition~\ref{prop:Mw0}). These are proved here, given the
  existence of normalizing factors, Theorem~\ref{thm:normalization}, which
  is taken from~\cite{Ar89,Li12a}.
\item \emph{The stable(r) form} (Section~\ref{sec:stable}). Conjugation
  by $GL(2)$ can exchange the two lifts of an element
  (Lemma~\ref{lem:eta}). Stable conjugacy is conjugation by $GL(2)$
  twisted by the sign $(\det h,\det(1+\gamma))_v$, and its classes are
  described by the function $\Delta_\psi$ (Lemma~\ref{lem:Delta}). The
  character of $r^+_\psi-r^-_\psi$ is computed on the regular set
  (Proposition~\ref{prop:weil-char}). The regular elliptic terms take the
  form $S_{\mathrm{ell}}(f)+S_{\mathrm{ell}}(f_-)$
  (Theorem~\ref{thm:elliptic-stable}). The unipotent coefficients are
  $a^G(S,u_\xi)=\frac12\bigl(c_S+\sum_{\kappa\ne1}\kappa_S(\xi)
  L^S(1,\kappa)\bigr)$ (Proposition~\ref{prop:unipotent}). The theta
  functions contribute to the spectral side in the same form as the
  elliptic terms (Proposition~\ref{prop:theta-packet}). The results are
  collected in Theorem~\ref{thm:stabler}. Its statements about the
  weighted terms use the property~\ref{N5} of the normalizing factors.
  Factors with this property exist at every place
  (Proposition~\normref{}(a),~(b)). Their explicit form at the dyadic and
  ramified places, Proposition~\normref{}(c), is taken from
  Szpruch~\cite{Sz09,Sz16}.
\end{enumerate}

\subsection*{Open questions}

The following questions are named as open in Sections~\ref{sec:intro}
and~\ref{sec:stable}. They are not answered here.

\begin{enumerate}[label=\arabic*.]
\item The stabilization of the weighted terms $I_A(a(t),\cdot)$ and
  $I_A(z,\cdot)$, which are equivariant under $GL(2)$ but not stable
  (Proposition~\ref{prop:IA-gl2} and Section~\ref{sec:stable-form}).
\item The transfer $f\mapsto f'$ to $SO(3)$ and the corresponding
  fundamental lemma~\cite{Li11,Li12b}, and with them the identification
  of $\Delta_\psi$ with Li's transfer factor and of $J^{\mathrm{st}}_G$
  with his stable orbital integral (Remark~\ref{rem:sec7-status}).
\item The description of the cuspidal spectrum by
  packets~\cite{Wa80,Wa91,GI18}. The character
  identity~\eqref{eq:packet-lift} used for the packets is quoted, not
  proved.
\item Whether the normalizing factors of~\cite{Li12a} have the
  property~\ref{N5}, and an independent derivation of Szpruch's formula
  for the local coefficients at the dyadic places
  (Remark~\ref{rem:N5}).
\item Whether $T\ge0$ suffices in Theorem~\ref{thm:Jchi-poly}, as it does
  on the geometric side (Remark~\ref{rem:sec5-status}).
\item The limit formulas of Section~\ref{sec:unipotent} at the singular
  elliptic points.
\item The $n$-fold cover for $n>2$, which the 1995 notes treated and
  this revision does not (Remarks~\ref{rem:general-n}
  and~\ref{rem:general-n-rank0}).
\end{enumerate}
